% ==================================================================
% id:       cond-028
% topic:    הסתברות מותנית ואי-תלות
% subtopic: מאורעות בלתי תלויים
% solution: full
% ==================================================================
$60\%$ מהסטודנטים בקורס בהסתברות גרים בבאר שבע. $40\%$ מהגרים בבאר שבע לומדים בקבוצת התרגול הראשונה, והשאר בקבוצה השנייה. $80\%$ מהלומדים בקבוצה הראשונה גרים בבאר שבע. בוחרים סטודנט מהקורס באקראי.
\begin{enumerate}
\item מה ההסתברות שהוא לומד בקבוצה הראשונה?
\begin{solbox}
נסמן $R$ = ״גר בבאר שבע״, $G$ = ״לומד בקבוצה הראשונה״. הנתונים: $\PP(R)=0.6$, $\PP(G\mid R)=0.4$, $\PP(R\mid G)=0.8$. לפי כלל הכפל
\[
\PP(R\cap G)=\PP(G\mid R)\,\PP(R)=0.4\cdot0.6=0.24.
\]
מצד שני, $\PP(R\cap G)=\PP(R\mid G)\,\PP(G)=0.8\,\PP(G)$, ולכן $\PP(G)=\frac{0.24}{0.8}=\ans{0.3}$.
\end{solbox}
\item האם המאורעות ״הסטודנט גר בבאר שבע״ ו־״הסטודנט לומד בקבוצה הראשונה״ בלתי תלויים?
\begin{solbox}
$\PP(R)\,\PP(G)=0.6\cdot0.3=0.18\neq0.24=\PP(R\cap G)$, ולכן המאורעות \(\ans{\text{תלויים}}\). (אפשר לראות את זה גם כך: $\PP(R\mid G)=0.8\neq0.6=\PP(R)$.)
\end{solbox}
\end{enumerate}
